\chapter*{Resumen}
\label{sec:Resumen}
\noindent La dinámica actual del mercado laboral, impulsada por la automatización y la inteligencia artificial, exige herramientas precisas para la normalización de competencias y ocupaciones laborales. En este contexto, el emparejamiento automático de títulos de ofertas de empleo se presenta como un desafío fundamental en el procesamiento del lenguaje natural.

Este trabajo describe el sistema desarrollado para la participación en la tarea de \glslink{JTM}{Job Title Matching} de la competición \glslink{Tarea}{TalentCLEF}, donde se pretende abordar el caso para el dominio del español. Se propone una arquitectura modular en cascada diseñada con la intención de maximizar la precisión en la recuperación de candidatos. El sistema integra tres etapas secuenciales: una fase de recuperación densa basada en tres modelos bi-encoders optimizados mediante aprendizaje supervisado y funciones de pérdida contrastivas; una etapa de fusión de rankings que combina las predicciones para mejorar la robustez; y una fase final de re-ranking mediante un modelo cross-encoder para refinar la selección de candidatos.

El desarrollo experimental incluyó un análisis exhaustivo de modelos en escenario \textit{zero-shot}, revelando la superioridad de las arquitecturas multilingües sobre las monolingües, incluso para el contexto del idioma español. Los resultados obtenidos validan la eficacia del enfoque propuesto: el sistema alcanzó la tercera posición por participantes en la clasificación general, con una métrica \gls{MAP} media de 0.521 en los tres idiomas obligatorios a presentar. Sin embargo, de los tres idiomas (inglés, alemán y español), el rendimiento fue más bajo en el idioma objetivo. El análisis posterior discute la discrepancia entre el rendimiento en validación y test, así como la paradoja de obtener mejores resultados en otros idiomas pese a la optimización específica para el español, subrayando la importancia de la diversidad de datos y los sesgos intrínsecos que caracterizan a los modelos preentrenados de forma masiva.